% Apuntes universitarios: Derecho de la Integración
% Mercosur como espacio de integración regional
% Método Cornell

\documentclass[12pt,oneside]{book}

% ===== METADATOS =====
\newcommand{\universidad}{Universidad de Buenos Aires}
\newcommand{\facultad}{Facultad de Derecho}
\newcommand{\carrera}{}
\newcommand{\materia}{Derecho de la Integración}
\newcommand{\titulo}{El Mercosur como espacio de integración regional}
\newcommand{\autor}{Isaac Edgar Camacho Ocampo}
\newcommand{\anio}{2026}

% ===== PAQUETES =====
\usepackage[utf8]{inputenc}
\usepackage[T1]{fontenc}
\usepackage[spanish,es-nodecimaldot]{babel}

\usepackage[
    top=2.5cm,
    bottom=2.5cm,
    left=3cm,
    right=2.5cm,
    headheight=15pt
]{geometry}

\usepackage{amsmath}
\usepackage{amssymb}
\usepackage{mathtools}

\usepackage{graphicx}
\usepackage{float}
\usepackage{caption}
\usepackage{subcaption}

\usepackage{booktabs}
\usepackage{array}
\usepackage{longtable}
\usepackage{tabularx}

\usepackage{enumitem}

\usepackage{xcolor}
\usepackage[most]{tcolorbox}

\usepackage{fancyhdr}

\usepackage[
    colorlinks=true,
    linkcolor=blue!50!black,
    citecolor=green!40!black,
    urlcolor=blue!60!black,
    pdftitle={\titulo},
    pdfauthor={\autor},
    pdfsubject={Apuntes universitarios},
    bookmarks=true
]{hyperref}

% ===== CONFIGURACIÓN DE PÁRRAFOS =====
\setlength{\parindent}{0pt}
\setlength{\parskip}{0.5em}

% ===== CONFIGURACIÓN DE LISTAS =====
\setlist[itemize]{
    topsep=0.3em,
    itemsep=0.2em
}

\setlist[enumerate]{
    topsep=0.3em,
    itemsep=0.2em
}

% ===== CONTROL DE VIUDAS Y HUÉRFANAS =====
\clubpenalty=10000
\widowpenalty=10000

% ===== ENCABEZADOS Y PIES =====
\pagestyle{fancy}
\fancyhf{}

\fancyhead[L]{\nouppercase{\leftmark}}
\fancyhead[R]{\materia}

\fancyfoot[C]{\thepage}

\renewcommand{\headrulewidth}{0.4pt}
\renewcommand{\footrulewidth}{0pt}

\fancypagestyle{plain}{
    \fancyhf{}
    \fancyfoot[C]{\thepage}
    \renewcommand{\headrulewidth}{0pt}
}

% ===== COMANDOS MATEMÁTICOS =====
\newcommand{\abs}[1]{\left|#1\right|}
\newcommand{\norm}[1]{\left\lVert#1\right\rVert}
\newcommand{\vect}[1]{\mathbf{#1}}
\newcommand{\deriv}[2]{\frac{d#1}{d#2}}
\newcommand{\pderiv}[2]{\frac{\partial#1}{\partial#2}}

% ===== DEFINICIÓN DE CAJAS ACADÉMICAS =====

\newtcolorbox{definition}[1]{
    enhanced,
    breakable,
    title={Definición: #1},
    fonttitle=\bfseries,
    colback=blue!3,
    colframe=blue!50!black,
    boxrule=0.8pt,
    arc=2mm
}

\newtcolorbox{important}{
    enhanced,
    breakable,
    title={Importante},
    fonttitle=\bfseries,
    colback=yellow!8,
    colframe=orange!70!black,
    boxrule=0.8pt,
    arc=2mm
}

\newtcolorbox{example}[1]{
    enhanced,
    breakable,
    title={Ejemplo: #1},
    fonttitle=\bfseries,
    colback=green!4,
    colframe=green!40!black,
    boxrule=0.8pt,
    arc=2mm
}

\newtcolorbox{note}{
    enhanced,
    breakable,
    title={Nota},
    fonttitle=\bfseries,
    colback=white,
    colframe=gray!50,
    boxrule=0.6pt,
    arc=2mm
}

% ===== ÍNDICE =====
\setcounter{tocdepth}{2}

% ===== INICIO DEL DOCUMENTO =====
\begin{document}

% ===== PORTADA =====
\begin{titlepage}

\thispagestyle{empty}

\begin{center}

\vspace*{1cm}

{\large \textsc{\universidad}}

\vspace{0.2cm}

{\large \textsc{\facultad}}

\vspace{3cm}

{\Huge \textbf{\materia}}

\vspace{1cm}

{\LARGE \titulo}

\vspace{0.8cm}

\textit{La integración regional es el primer paso para transformar la vecindad en comunidad.}

\vspace{1.2cm}

\rule{0.70\textwidth}{0.5pt}

\vspace{0.5cm}

{\large Apuntes de estudio}

\vspace{0.5cm}

\rule{0.70\textwidth}{0.5pt}

\vfill

{\large \autor}

\vspace{0.3cm}

{\large \anio}

\end{center}

\vfill

\begin{flushleft}
\textit{Este es un modesto aporte a los estudiantes de ``Derecho de la Integración'' no pretende sustituir el material oficial de la materia.}
\end{flushleft}

\end{titlepage}

% ===== ÍNDICE =====
\tableofcontents
\newpage

% ===== CAPÍTULO 1: ANTECEDENTES EN LA ALADI =====
\chapter{Antecedentes del Mercosur en el marco de la ALADI}

\section{La ALADI como marco de integración}

El Mercosur no constituye un espacio de integración autónomo o independiente, sino que se inscribe dentro de un marco institucional más amplio: la Asociación Latinoamericana de Integración (ALADI). Esta distinción resulta fundamental para comprender la naturaleza del proceso de integración regional en América Latina.

La ALADI, en su estructura normativa, permite a sus Estados miembros celebrar acuerdos de alcance parcial. Estos acuerdos, como su nombre indica, no comprometen la participación de la totalidad de los miembros del organismo, sino que vinculan únicamente a los Estados que los suscriben. De esta manera, la ALADI autoriza, al mismo tiempo, la conformación de espacios de integración subregionales entre algunos de sus miembros.

\begin{definition}{Acuerdo de alcance parcial}
Instrumento jurídico que permite a algunos miembros de la ALADI celebrar acuerdos específicos de integración sin que estos comprometan a la totalidad de los Estados miembros del organismo ni contradigan el objetivo general de integración más amplio.
\end{definition}

El Mercosur, en consecuencia, es uno de estos espacios subregionales: no representa un desprendimiento de la ALADI ni algo distinto de ella. Por el contrario, se trata de un acuerdo que se da \textbf{dentro} de la ALADI y que, lejos de contradecir los objetivos generales del organismo, apunta a consolidar relaciones cada vez más estrechas entre algunos de sus miembros. Esta lógica permite que otros miembros de la ALADI puedan sumarse posteriormente al acuerdo sin que esto implique alterar el marco general.

\section{Antecedentes históricos: el acercamiento entre Argentina y Brasil}

El proceso que desembocó en la creación del Mercosur se remonta al restablecimiento de la democracia en los países de América Latina durante la década de 1980. En ese contexto, Argentina y Brasil —bajo la presidencia de Raúl Alfonsín en el caso argentino— comenzaron un proceso de aproximación a través de negociaciones bilaterales desarrolladas dentro del ámbito de la ALADI.

Este acercamiento entre dos de los mayores actores económicos y políticos de la región fue gradual. Durante este período se celebraron numerosos acuerdos de alcance parcial que sentaron las bases para una relación cada vez más profunda. Estos acuerdos previos funcionaron como un ensayo de lo que vendría después: permitieron a ambos países explorar campos de cooperación, ajustar expectativas y generar confianza mutua.

\section{Formación del bloque: los Estados fundadores}

Continuando ese proceso de acercamiento, ya consolidada la democracia en la región, Argentina, Brasil, Paraguay y Uruguay decidieron avanzar hacia un acuerdo de integración de mayor alcance y profundidad. Estos cuatro países son, precisamente, los \textbf{Estados fundadores del Mercosur}.

\begin{important}
La decisión de formar el Mercosur fue, fundamentalmente, política: representó una apuesta por consolidar la integración regional como proyecto estratégico de largo plazo, más allá de las variaciones en las coyunturas económicas o en los gobiernos de turno.
\end{important}

% ===== CAPÍTULO 2: EL TRATADO DE ASUNCIÓN =====
\chapter{El Tratado de Asunción y la creación del Mercado Común del Sur}

\section{Naturaleza y alcance del tratado}

\begin{definition}{Tratado de Asunción}
Instrumento jurídico internacional suscripto en 1991 por Argentina, Brasil, Paraguay y Uruguay, que crea el Mercado Común del Sur (Mercosur) y constituye el derecho originario del espacio de integración.
\end{definition}

El Tratado de Asunción cumple, dentro del Mercosur, una función equivalente a la que desempeña una Constitución dentro del ordenamiento jurídico interno de un Estado. Esta equivalencia funcional resulta decisiva para entender la jerarquía normativa del bloque y para distinguir entre las normas que conforman su estructura fundamental y las que derivan de esa estructura.

\section{El concepto de mercado común}

Un mercado común, como estadio de integración económica, aspira a lograr la \textbf{libre circulación de mercaderías, servicios y personas} entre los Estados que lo componen. Además, busca garantizar la \textbf{libre circulación de los factores de producción}: el capital y el trabajo.

Estas libertades fundamentales requieren, como requisito previo, la adopción de \textbf{políticas uniformes} entre los Estados miembros. El Mercosur, en su formulación, establece este objetivo: constituir un mercado único cuyos participantes —empresas, trabajadores, inversores— puedan moverse sin las restricciones que normalmente imponen las fronteras nacionales.

Sin embargo, la realidad del bloque ha mostrado que, en materia de políticas uniformes, existe una considerable distancia entre el objetivo declarado y lo efectivamente alcanzado. El Mercosur debe evaluar, tema por tema, materia por materia, en qué medida ha logrado armonizar sus políticas y en qué medida persisten divergencias que limitan la concreción del ideal del mercado común.

\section{El costo de la integración: la cesión de facultades soberanas}

Formar parte de un mercado común implica necesariamente una cesión de facultades: los Estados \textbf{no pueden tomar decisiones unilaterales} en materias que afecten el libre tránsito de personas, capitales, servicios o mercaderías dentro del bloque. Cualquier decisión interna que restrinja estas libertades impacta directamente sobre el funcionamiento del mercado común y, por lo tanto, afecta a los demás Estados miembros.

Esta restricción de la soberanía decisoria constituye, precisamente, el núcleo de todo proceso de integración: los Estados aceptan límites a su capacidad de actuar autónomamente para lograr, a cambio, los beneficios que la integración ofrece.

\subsection{Tensión entre política interna y compromisos internacionales}

Una de las dificultades recurrentes en el funcionamiento del Mercosur surge de la tensión entre las decisiones que adopta cada Estado en su ámbito interno y los compromisos que ha asumido frente a sus socios del bloque. Frecuentemente, los responsables políticos de un Estado toman decisiones atendiendo a la coyuntura inmediata, sin advertir que esas decisiones —aunque justificadas por circunstancias internas— pueden resultar incompatibles con los compromisos externos previamente asumidos.

\begin{example}{El voto \textit{no positivo} de 2008}
Durante un conflicto relativo a la regulación de retenciones a las exportaciones del sector agropecuario, el Senado de la Nación Argentina se encontró en una situación de empate al momento de votar un proyecto de ley. La presidencia del Senado debió ejercer su voto de desempate, que resultó siendo negativo, rechazando así la iniciativa.

Lo particularmente significativo del episodio radica en su timing: la misma semana en que el Senado argentino rechazaba esta iniciativa de apertura comercial, el gobierno nacional estaba suscribiendo, en el ámbito del Mercosur, compromisos de libre circulación con sus socios regionales.

Esta incoherencia ejemplifica la contradicción frecuente: un Estado limita la libre circulación de mercaderías en su territorio para proteger un sector productivo interno, pero simultáneamente se compromete internacionalmente a permitir esa libre circulación. Cualquier límite que un Estado imponga debería dirigirse únicamente a terceros países ajenos al bloque; nunca debe dirigirse contra sus propios socios, porque de lo contrario se estaría violando directamente el compromiso asumido en el Mercosur.
\end{example}

Esta clase de incoherencias es frecuente y genera tensiones permanentes dentro del proceso de integración. La razón subyacente es que los gobiernos enfrentan presiones políticas internas que, en el corto plazo, generan incentivos para adoptar medidas proteccionistas, mientras que los beneficios de la integración tienden a manifestarse a largo plazo y, además, a distribuirse de manera desigual entre los distintos sectores de la economía.

% ===== CAPÍTULO 3: DERECHO ORIGINARIO Y DERECHO DERIVADO =====
\chapter{Fuentes del Derecho de la Integración}

\section{Introducción a las fuentes}

Una de las distinciones más centrales en el derecho de la integración —y no exclusiva del Mercosur, sino aplicable a otros espacios de integración regional— es la que existe entre dos categorías fundamentales de normas: el \textbf{derecho originario} y el \textbf{derecho derivado}.

Esta distinción no es simplemente académica: tiene implicaciones prácticas decisivas. Determinar si una norma pertenece a una u otra categoría permite establecer su posición en la jerarquía normativa del bloque, qué autoridad la dictó, y qué procedimiento se requiere para modificarla.

\section{Derecho originario}

\begin{definition}{Derecho originario}
Conjunto de normas de fuente convencional que conforman la estructura fundamental del espacio de integración. En el Mercosur, está constituido por el Tratado de Asunción y por aquellos protocolos posteriores que, aunque sean instrumentos jurídicos distintos, se incorporan al tratado con su mismo rango y lo modifican estructuralmente.
\end{definition}

El derecho originario es, dentro del sistema de integración, lo que la Constitución es dentro del ordenamiento jurídico de un Estado: la norma fundamental que establece la estructura, los órganos y las reglas básicas del sistema.

Una confusión frecuente surge de la distinción entre tener \textbf{rango} constitucional y ser \textbf{parte} de la constitución. El Protocolo de Ouro Preto, que modificó la estructura orgánica del Mercosur, no simplemente \textbf{tiene} rango equivalente al Tratado de Asunción: pasa a \textbf{integrar}, junto con él, el propio derecho originario del sistema. No son dos normas de igual jerarquía; es una única fuente normativa fundamental en la que ambos instrumentos se han fusionado.

\subsection{El Protocolo de Ouro Preto}

El Protocolo de Ouro Preto es el instrumento que modificó la estructura institucional del Mercosur. Aunque se trata de un protocolo —es decir, un instrumento jurídico distinto del tratado original—, su incorporación al derecho originario fue total: pasó a formar parte integral de la norma fundacional del bloque, alterando su estructura orgánica y sus mecanismos de funcionamiento.

\section{Derecho derivado}

\begin{definition}{Derecho derivado}
Conjunto de normas dictadas por los órganos que el tratado fundacional creó, en ejercicio de las funciones que ese tratado les atribuyó. Incluye resoluciones, decisiones, directivas y otros actos normativos emanados de la estructura institucional del Mercosur.
\end{definition}

De manera análoga a como la Constitución Nacional crea los poderes Ejecutivo, Legislativo y Judicial, el Tratado de Asunción crea una estructura orgánica y dota a cada órgano de determinadas funciones y competencias. La actividad de esos órganos —cuando actúan conforme a las competencias que el tratado les asignó— constituye derecho derivado.

\subsection{Criterio de distinción}

El criterio fundamental para distinguir entre derecho originario y derecho derivado es la \textbf{cualidad en que actúa} la autoridad emisora:

\begin{itemize}
\item Cuando los presidentes de los Estados miembros se reúnen y firman un instrumento \textbf{actuando en representación de su Estado} (como ocurrió con el Tratado de Asunción), están generando derecho originario: una norma de fuente convencional internacional.

\item Cuando esos mismos presidentes —u otros funcionarios— actúan integrando un órgano creado por el tratado (por ejemplo, el Consejo del Mercado Común), lo que producen es derecho derivado, aunque se trate de las mismas personas físicas en su capacidad individual.
\end{itemize}

La persona que firma puede ser la misma, pero la función que desempeña en ese momento es decisiva para clasificar la norma que produce.

\section{Tabla comparativa: derecho originario y derivado}

\begin{table}[H]
\centering
\begin{tabularx}{\textwidth}{
    >{\raggedright\arraybackslash}p{0.25\textwidth}
    >{\raggedright\arraybackslash}X
    >{\raggedright\arraybackslash}X
}
\toprule
\textbf{Aspecto} & \textbf{Derecho originario} & \textbf{Derecho derivado} \\
\midrule
Equivalente interno & Constitución Nacional & Leyes, resoluciones, sentencias \\
Instrumento & Tratado de Asunción y protocolos que lo modifican (ej. Ouro Preto) & Resoluciones, decisiones y directivas de los órganos del Mercosur \\
Quién lo dicta & Los Estados Parte, a través de sus presidentes, en representación de su Estado & Los órganos creados por el tratado, actuando dentro de su competencia \\
Función & Crea el sistema, la estructura orgánica y las reglas fundamentales & Desarrolla y ejecuta las funciones atribuidas por el derecho originario \\
Modificación & Requiere nueva convención de los Estados & Puede modificarse dentro de los límites de la competencia del órgano \\
\bottomrule
\end{tabularx}
\caption{Distinción entre derecho originario y derivado en el Mercosur}
\end{table}

\section{Implicaciones prácticas}

Para el ejercicio profesional, la distinción entre ambas fuentes es fundamental. Determinar si una norma es derecho originario o derivado, y verificar que el órgano que la dictó era competente para hacerlo, es la manera de establecer si esa norma resulta exigible. Si el órgano no tenía competencia para dictar esa norma, entonces la norma no existe jurídicamente, aunque formalmente haya sido emitida, y no genera obligación de acatamiento.

\begin{important}
Una resolución de un órgano del Mercosur que exceda sus competencias no obliga a los Estados miembros. La verificación de competencia es, por lo tanto, parte indispensable del análisis de cualquier acto normativo del bloque.
\end{important}

% ===== CAPÍTULO 4: ESTRUCTURA INSTITUCIONAL =====
\chapter{Estructura institucional del Mercosur}

\section{Introducción a los órganos}

El Tratado de Asunción crea una estructura institucional compuesta por tres órganos con capacidad decisoria. Estos órganos, que serán analizados en orden jerárquico, presentan distintos niveles de composición y funciones específicas dentro del sistema de integración.

\section{Consejo del Mercado Común (CMC)}

\begin{definition}{Consejo del Mercado Común}
Órgano de máxima autoridad y nivel decisorio más alto del Mercosur, responsable de definir la política general y el rumbo estratégico del bloque.
\end{definition}

\subsection{Composición}

El Consejo del Mercado Común está integrado por:

\begin{itemize}
\item Los presidentes de los Estados miembros del Mercosur
\item Los ministros de Economía de cada Estado
\item Los cancilleres (ministros de Relaciones Exteriores)
\end{itemize}

La composición en este nivel máximo responde a la naturaleza de las decisiones que se toman: se trata de decisiones que afectan simultáneamente dimensiones económicas (de allí la participación de los ministros de Economía) y de política exterior (de allí la de los cancilleres), y que requieren, en última instancia, la aprobación política de cada jefe de Estado.

\subsection{Función}

El Consejo es el órgano que \textbf{``maneja el timón''} del Mercosur: traza el rumbo político del bloque. Define orientaciones estratégicas tales como:

\begin{itemize}
\item Si conviene priorizar la negociación con el mercado latinoamericano o con otros bloques regionales
\item La relación política general con otros espacios de integración (por ejemplo, la Unión Europea)
\item Cuestiones fundamentales de política comercial y aranceles
\item Aprobación de nuevos miembros o suspensiones
\end{itemize}

Las decisiones del Consejo establecen las líneas maestras que los demás órganos implementan.

\section{Grupo Mercado Común (GMC)}

\begin{definition}{Grupo Mercado Común}
Órgano ejecutivo e implementador de las políticas definidas por el Consejo del Mercado Común, responsable de poner en marcha las decisiones estratégicas del bloque.
\end{definition}

\subsection{Composición}

El GMC está integrado por funcionarios del Ministerio de Economía de cada Estado, aunque no necesariamente por los ministros personalmente, sino por equipos técnicos de nivel inferior. Esta estructura técnica refleja que el Grupo no es un órgano de toma de decisiones políticas de nivel máximo, sino de ejecución e implementación.

\subsection{Estructura de trabajo}

El GMC se organiza internamente a través de \textbf{subgrupos de trabajo} temáticos. Cada subgrupo se especializa en un área particular de la integración. A modo de ejemplo, la documentación de clase menciona la existencia de un \textbf{Subgrupo de Trabajo N.º 7} dedicado a cuestiones de protección del consumidor.

Estos subgrupos permiten una especialización temática que facilita la profundización en materias específicas.

\subsection{Función}

El rol del GMC es ejecutar y operacionalizar las decisiones estratégicas del Consejo. Si el Consejo decide que deben eliminarse ciertos aranceles entre los miembros, es el GMC quien diseña el cronograma, identifica los obstáculos prácticos, coordina con las autoridades nacionales, y propone los mecanismos específicos de implementación.

\section{Comisión de Comercio del Mercosur (CCM)}

\begin{definition}{Comisión de Comercio del Mercosur}
Órgano especializado responsable de la administración técnica de las cuestiones comerciales entre los Estados miembros, incluyendo regulación de aranceles y barreras comerciales.
\end{definition}

\subsection{Composición y estructura}

La CCM está integrada por funcionarios especializados en comercio de cada Estado. Al igual que el GMC, se organiza internamente a través de \textbf{subcomités organizados por materias}. A modo de ilustración, la clase menciona la existencia de un subcomité en materia agropecuaria.

\subsection{Función}

La CCM se ocupa específicamente de las cuestiones comerciales dentro del bloque:

\begin{itemize}
\item Administración de aranceles (derechos de importación y exportación)
\item Regulación de barreras arancelarias (impuestos y derechos)
\item Regulación de barreras para-arancelarias (normas técnicas, sanitarias, etc.)
\item Todas las cuestiones vinculadas con el intercambio de mercaderías dentro del bloque
\end{itemize}

\section{Tabla comparativa: estructura institucional}

\begin{table}[H]
\centering
\begin{tabularx}{\textwidth}{
    >{\raggedright\arraybackslash}p{0.22\textwidth}
    >{\raggedright\arraybackslash}X
    >{\raggedright\arraybackslash}X
}
\toprule
\textbf{Órgano} & \textbf{Composición} & \textbf{Función principal} \\
\midrule
Consejo del Mercado Común (CMC) & Presidentes, ministros de Economía y cancilleres & Define la política y el rumbo general del bloque \\
Grupo Mercado Común (GMC) & Funcionarios de Ministerios de Economía, organizados en subgrupos temáticos & Ejecuta e implementa las políticas fijadas por el Consejo \\
Comisión de Comercio del Mercosur (CCM) & Funcionarios especializados, organizados en subcomités temáticos & Regula y administra las cuestiones estrictamente comerciales \\
\bottomrule
\end{tabularx}
\caption{Estructura institucional del Mercosur}
\end{table}

\section{Naturaleza intergubernamental de los órganos}

Un aspecto fundamental de la estructura del Mercosur es que todos estos órganos —el Consejo, el Grupo y la Comisión— son órganos \textbf{intergubernamentales}. Esto significa que están integrados por representantes de los gobiernos de cada Estado miembro, y no por funcionarios independientes ni por representantes directos de la ciudadanía del bloque.

Esta característica marca una diferencia fundamental respecto de otros modelos de integración: los órganos del Mercosur no tienen la independencia respecto de los gobiernos nacionales que caracteriza a órganos análogos en otros espacios de integración (como la Comisión Europea).

% ===== CAPÍTULO 5: CONSENSO Y ESTATUS DE LOS ESTADOS =====
\chapter{Mecanismo de decisión: el consenso}

\section{El consenso como regla de decisión}

\begin{definition}{Consenso}
Mecanismo de adopción de normas que exige unanimidad: la norma se dicta únicamente si todos los órganos y miembros que participan en la decisión llegan a un acuerdo. Un solo voto en contra impide que la norma llegue a nacer.
\end{definition}

Los tres órganos del Mercosur con capacidad decisoria ---el Consejo, el Grupo Mercado Común y la Comisión de Comercio--- adoptan sus normas por consenso, es decir, por unanimidad.

La consecuencia de este mecanismo es decisiva: \textbf{si un solo Estado miembro no está de acuerdo}, la norma directamente no se dicta. No es que la norma exista pero no se aplique a ese Estado, sino que la norma \textbf{no llega a nacer}. No existe jurídicamente. No genera obligación para ningún Estado.

\begin{important}
El consenso es un mecanismo binario: o todos acuerdan y la norma existe, u al menos uno se opone y la norma simplemente no nace.
\end{important}

\section{Efecto práctico: la no obligación del disidente}

Un Estado que se opone a una norma no queda obligado por ella simplemente porque los demás la aprobaron. Este es un principio fundamental del derecho internacional convencional: nadie se obliga sin su consentimiento. El Mercosur, en su mecanismo de consenso, mantiene esta lógica.

Si el Mercosur decidiera por mayoría, entonces sería posible que una norma se aprobara sin el voto de todos los miembros, y todos quedarían obligados por ella. Pero con el sistema de consenso, cada Estado tiene poder de veto, lo que significa que cada Estado tiene la capacidad de impedir que se dicte cualquier norma a la que se oponga.

\section{Comparación con la Unión Europea}

\subsection{Sistema de mayorías en la Unión Europea}

La Unión Europea, en contraste directo, adopta las decisiones según distintos tipos de mayorías:

\begin{itemize}
\item Mayorías simples (para ciertos temas de menor importancia)
\item Mayorías absolutas (para asuntos de relevancia intermedia)
\item Mayorías calificadas (para decisiones de mayor trascendencia)
\item Dobles mayorías (en ciertos casos, requieren tanto mayoría de votos como representación de población)
\end{itemize}

El sistema de mayorías implica que una norma puede aprobarse aún cuando algunos Estados se oponen. Los Estados disidentes no pueden impedir la norma por voto de veto; deben acatarla igual, aunque no estén de acuerdo.

\subsection{Contexto histórico de la Unión Europea}

Cuando la Unión Europea (en su forma primitiva como Comunidad Económica Europea) fue fundada en 1957, estaba integrada por apenas seis Estados miembros. Incluso con tan pocos miembros, sus fundadores optaron por un sistema de mayorías en lugar de unanimidad, reconociendo que la unanimidad habría resultado impracticable.

A través de las décadas, la Unión Europea fue expandiéndose. Hoy cuenta con veintisiete Estados miembros. Con esa cantidad de miembros, una regla de unanimidad habría hecho prácticamente imposible dictar cualquier norma, porque siempre habría habido al menos un Estado en desacuerdo.

\subsection{Estructura institucional más compleja}

Además del mecanismo de decisión, la estructura institucional de la Unión Europea es notablemente más compleja que la del Mercosur. La Unión cuenta con órganos que representan distintos intereses y funciones:

\begin{itemize}
\item Órganos que representan a la propia Unión como entidad supranacional (la Comisión Europea)
\item Órganos que representan a los Estados miembros (el Consejo de la Unión)
\item Órganos que representan a la población (el Parlamento Europeo)
\end{itemize}

Esta multiplicidad de órganos con representaciones distintas constituye una forma de división de poderes ausente en el esquema del Mercosur, donde todos los órganos son puramente intergubernamentales.

\section{Tabla comparativa: Mercosur versus Unión Europea}

\begin{table}[H]
\centering
\begin{tabularx}{\textwidth}{
    >{\raggedright\arraybackslash}p{0.2\textwidth}
    >{\raggedright\arraybackslash}X
    >{\raggedright\arraybackslash}X
}
\toprule
\textbf{Aspecto} & \textbf{Mercosur} & \textbf{Unión Europea} \\
\midrule
Mecanismo de decisión & Consenso (unanimidad) & Mayoría (simple, absoluta, calificada, doble) \\
Efecto de un voto en contra & La norma no se dicta & La norma puede aprobarse igual, sin ese Estado \\
Representación en los órganos & Solo representantes de los gobiernos (intergubernamental) & Representan a la Unión, a los Estados y a la población \\
Lógica dominante & Negociación internacional & Integración supranacional \\
Cantidad de miembros originarios & 4 (Argentina, Brasil, Paraguay, Uruguay) & 6 (Francia, Alemania, Italia, Bélgica, Holanda, Luxemburgo) \\
\bottomrule
\end{tabularx}
\caption{Comparación entre los sistemas de decisión y estructuras del Mercosur y la Unión Europea}
\end{table}

\section{Significado político del consenso}

La exigencia de consenso revela algo profundo sobre la naturaleza política del Mercosur: \textbf{los Estados miembros no están dispuestos a someterse a la voluntad de la mayoría}. Cada Estado se reserva el derecho de veto sobre cualquier decisión que considere que afecta sus intereses.

Esta característica evidencia, además, la \textbf{ausencia de una voluntad política real de avanzar hacia órganos supranacionales}. Si existiera esa voluntad, la estructura del bloque sería distinta: habría órganos independientes con capacidad de decisión, con el reconocimiento de que ciertos asuntos quedan fuera del control exclusivo de cada Estado.

La persistencia del consenso muestra que, en última instancia, cada Estado prioriza su propia posición frente a la del bloque en conjunto.

% ===== CAPÍTULO 6: ESTATUS DE LOS ESTADOS Y SUSPENSIONES =====
\chapter{Socios plenos, Estados asociados y suspensiones}

\section{Categorías de participación en el Mercosur}

El Mercosur distingue entre dos categorías de Estados con grados distintos de vinculación y derechos:

\begin{enumerate}
\item Estados socios plenos (o simplemente \textbf{Estados parte})
\item Estados asociados
\end{enumerate}

\section{Estados socios plenos}

\begin{definition}{Estado socio pleno}
Estado que se ha incorporado plenamente al Mercosur, asumiendo la obligación de implementar la totalidad de la normativa del bloque y gozando del derecho a participar en todas las decisiones a través del voto.
\end{definition}

\subsection{Características}

\begin{itemize}
\item \textbf{Alcance de la obligación}: El Estado está obligado por la \textit{totalidad} de la normativa del bloque, tanto originaria como derivada.
\item \textbf{Derecho a voto}: Goza de plenos derechos de voto en los órganos con capacidad decisoria (CMC, GMC, CCM).
\item \textbf{Participación}: Participa plenamente como miembro del bloque en todas sus dimensiones.
\end{itemize}

\subsection{Estados fundadores y posteriores}

Los cuatro Estados fundadores del Mercosur son:

\begin{itemize}
\item Argentina
\item Brasil
\item Paraguay
\item Uruguay
\end{itemize}

Posteriormente, \textbf{Venezuela} se incorporó como socio pleno, aunque posteriormente ha sido suspendida del bloque.

\section{Estados asociados}

\begin{definition}{Estado asociado}
Estado que mantiene una relación de integración parcial con el Mercosur, obligándose únicamente por la normativa a la que adhiere expresamente, sin derecho a participar en la toma de decisiones.
\end{definition}

\subsection{Características}

\begin{itemize}
\item \textbf{Alcance de la obligación}: El Estado asociado solo se obliga por la normativa del bloque a la que adhiere expresamente, no por la totalidad de las normas.
\item \textbf{Derecho a voto}: No tiene derecho a voto en los órganos decisorios.
\item \textbf{Participación}: Puede participar de las reuniones y estar presente, pero sin capacidad de decisión.
\end{itemize}

\subsection{Ejemplos de Estados asociados}

Entre los Estados asociados se encuentran:

\begin{itemize}
\item Chile
\item Colombia
\item Bolivia (que además mantiene un trámite de incorporación como socio pleno)
\end{itemize}

\subsection{Razón de la asociación parcial}

La condición de Estado asociado está ligada a la pertenencia de estos países a la ALADI. Por formar parte de ese marco más amplio, se les permite vincularse parcialmente con el Mercosur sin asumir la totalidad de sus compromisos, lo que genera una situación de integración gradual.

\section{Ingreso de Venezuela como socio pleno}

\subsection{Contexto político}

Venezuela se incorporó como socio pleno del Mercosur durante la presidencia de Hugo Chávez, en un contexto de cercanía política con el gobierno argentino de la época.

\subsection{Análisis económico y político}

Desde una perspectiva puramente económica, el ingreso de Venezuela no representó un aporte significativo para el bloque: la economía venezolana, aunque importante, no contribuía de manera determinante a fortalecer las economías del Mercosur ni ampliaba significativamente el mercado común.

Sin embargo, desde una perspectiva política, el ingreso de Venezuela cumplió una función estratégica: permitió conformar, frente a un Brasil cada vez más fuerte económicamente, un frente regional que compartiera una determinada orientación política y económica.

\section{Cláusula democrática y suspensiones}

\subsection{Protocolo de Ushuaia: la cláusula democrática}

Una de las premisas fundamentales del Mercosur es que sus Estados miembros deben operarse dentro del \textbf{Estado de Derecho} y en un ámbito \textbf{democrático}. Esta premisa se formaliza en la llamada \textbf{cláusula democrática del Mercosur} (referida en la documentación como ``Protocolo de Ushuaia'').

La cláusula democrática establece que aquellos Estados que se apartan de estos principios fundamentales de democracia y Estado de Derecho pueden ser suspendidos del bloque.

\subsection{Efectos de la suspensión}

Un Estado suspendido:

\begin{itemize}
\item Pierde el derecho a participar en el proceso de toma de decisiones (derecho a voto)
\item No puede gozar de los beneficios que corresponden a los Estados parte del Mercosur
\item Puede participar de las reuniones, pero sin capacidad decisoria
\item Permanece en esa condición hasta que se resuelva su reincorporación
\end{itemize}

\begin{important}
La suspensión es una medida que afecta tanto los derechos como las obligaciones del Estado suspendido: pierda la capacidad de decidir, pero también quedan limitados sus derechos de acceso al mercado común.
\end{important}

\section{El caso de Paraguay: suspensión e ingreso de Venezuela}

\subsection{Hechos del caso}

En un momento crítico de la historia del Mercosur, el mismo día en que se aprobó el ingreso de Venezuela como socio pleno, se había decidido previamente suspender a Paraguay del bloque. El gobierno paraguayo de ese momento era cuestionado por no ajustarse a las premisas democráticas establecidas en el Protocolo de Ushuaia.

\subsection{El rol de Paraguay en la decisión}

Paraguay era, precisamente, el Estado que habría votado en contra del ingreso de Venezuela, como parte de su capacidad de veto dentro del sistema de consenso. Al quedar suspendido, Paraguay perdió su derecho a voto, lo que permitió que el ingreso de Venezuela se aprobara sin que todos los miembros plenos completaran el requisito de unanimidad.

\subsection{Cuestión de validez posterior}

Cuando posteriormente se restableció el ingreso de Paraguay al bloque, ese país cuestionó la validez del ingreso de Venezuela, alegando que ``los miembros permanentes no votaron todos'' en esa sesión de aprobación. El argumento de Paraguay era que, sin su voto, no se había respetado el requisito de consenso.

\subsection{Interpretación política}

La clase caracteriza este episodio como un claro ejemplo de ``rosca política'': el uso de mecanismos formales (la cláusula democrática, la suspensión) para alterar el resultado de un proceso de toma de decisiones que, de otro modo, habría sido distinto.

\subsection{Situación actual}

En la actualidad, el Mercosur cuenta formalmente con cinco Estados socios plenos, incluida Venezuela. Sin embargo, Venezuela se encuentra a su vez suspendida del bloque, lo que genera una situación paradójica: es miembro pleno, pero está privada de los derechos y beneficios correspondientes a esa condición.

% ===== CAPÍTULO 7: TENSIONES ENTRE POLÍTICAS INTERNAS E INTEGRACIÓN =====
\chapter{Tensiones entre políticas internas y compromisos de integración}

\section{El desafío de sostener políticas de largo plazo}

Una de las dificultades más profundas que enfrenta el proceso de integración en el Mercosur no es ni institucional ni legal, sino política: la dificultad de que los gobiernos sostengan políticas de largo plazo, más allá de sus ciclos electorales.

Los gobiernos enfrentan presiones inmediatas: ciclos electorales cortos, demandas de sectores productivos, conflictos laborales. Estas presiones incentivan decisiones que privilegian lo inmediato sobre lo de largo plazo, aunque esas decisiones perjudiquen objetivos estratégicos que el propio Estado ha formalizado a nivel internacional.

\section{Ejemplo histórico: privatización de ferrocarriles y desindustrialización}

\subsection{Contexto}

Durante ciertos períodos de la historia económica argentina reciente, se tomaron decisiones de política económica que tuvieron consecuencias profundas a largo plazo. Una de las más significativas fue la privatización de los ferrocarriles del Estado.

\subsection{Efecto sobre el transporte de mercaderías}

La privatización de los ferrocarriles trasladó, en la práctica, el transporte de mercaderías desde el ferrocarril hacia el transporte automotor (camiones). Esta reconfiguración de la logística de transporte parece, en el corto plazo, una cuestión técnica menor; en realidad, reorganizó toda la estructura de costos y eficiencia de la economía.

\subsection{Impacto sobre la industria nacional}

De manera contemporánea, se produjo una apertura a productos importados que afectó gravemente a la industria nacional. Ejemplos concretos abundan: en barrios de la Ciudad de Buenos Aires como Parque Patricios, que históricamente albergaban numerosas fábricas de producción nacional (entre ellas una reconocida fábrica de chocolates), esas instalaciones fueron clausurando o reconvirtiéndose a otros usos.

La lógica combinada fue: se privatizó, se abrió a la competencia externa, y se redujo la protección a la industria nacional. El resultado fue que ``se cerraron'' las fábricas: sin protección de corto plazo, desaparecieron.

\subsection{Caracterización de la decisión}

Desde una perspectiva de largo plazo, estas decisiones destruyeron \textbf{capacidad productiva nacional}. Desde una perspectiva de corto plazo, pueden haber parecido beneficiosas (reducción de costos, acceso a productos más baratos). Pero la consecuencia fue la desaparición de sectores productivos que, una vez cerrados, resultan extremadamente difíciles de reconstruir.

\section{Consecuencias concretas: la escasez de neumáticos durante la pandemia}

\subsection{El episodio}

Durante la pandemia de COVID-19, la Argentina experimentó una escasez de neumáticos. La razón era que la capacidad productiva nacional de neumáticos había sido prácticamente eliminada durante los años en que se desmanteló la industria nacional.

\subsection{Impacto real}

La escasez de neumáticos no era un problema abstracto o de lujo. Los camiones son las arterias de la economía: sin neumáticos, los camiones no pueden circular. Si los camiones no circulan, no pueden transportar productos desde el campo hacia las ciudades o hacia los centros urbanos. El resultado es un problema concreto de abastecimiento de alimentos y productos esenciales para la población.

\subsection{Lección sobre decisiones de largo plazo}

Este episodio ilustra con crudeza cómo una decisión de política económica que parece razonable en el corto plazo (``dejamos que entren neumáticos importados más baratos'') genera consecuencias que, décadas después, resultan en crisis de abastecimiento.

\section{La necesidad de una política de Estado}

\subsection{La analogía del matrimonio}

Para ilustrar por qué resulta tan importante pensar en el largo plazo, la clase recurre a una analogía cotidiana:

\begin{quote}
``Cuando decidís casarte, no tenés que pensar en la noche de bodas; tenés que pensar en veinte años, cuando tengas que mandar a la universidad a tu hijo. (...) Tenés que pensar en las consecuencias que esto va a tener si estás haciendo un proyecto a largo plazo.''
\end{quote}

La misma lógica aplica a la política de integración: sostiene que sería deseable que los distintos gobiernos, más allá de su signo político, acordaran un conjunto mínimo de políticas de Estado que no se modifiquen con cada cambio de gestión.

\subsection{La falta de continuidad en las políticas}

Cuando cada gobierno que asume revierte lo hecho por el anterior, ningún proyecto de largo plazo llega a concretarse. Una industria requiere años para desarrollarse. La investigación requiere continuidad. Las políticas de integración requieren predictibilidad.

Si un gobierno protege cierto sector industrial con aranceles, pero el siguiente reduce esos aranceles, la industria no tiene estabilidad suficiente para invertir en tecnología, capacitación o expansión.

\subsection{Necesidad de consenso político transversal}

La clase propone que sería deseable que las distintas fuerzas políticas del país lograran acordar un conjunto de políticas de largo plazo que trascendieran los cambios de gobierno. No significa que todas las políticas sean fijas; significa que ciertos ejes estratégicos gocen de estabilidad más allá de los ciclos electorales.

\section{Ejemplo histórico: la guerra de Vietnam}

Como ilustración de la dificultad política de asumir los costos de ciertas decisiones, la clase menciona el caso de la guerra de Vietnam.

Durante casi cuarenta años, sucesivos presidentes de los Estados Unidos (de distintos partidos, con distintas ideologías) mantuvieron el compromiso con la guerra de Vietnam, a pesar de su costo humano y económico colosal. Ninguno de esos presidentes se atrevió, públicamente, a reconocer el error ni a ordenar un retiro que implicara asumir personalmente esa responsabilidad.

La razón no era técnica ni jurídica: era política. Asumir que se había cometido un error de esa magnitud habría tenido costos políticos enormes para quien lo reconociera.

Esta dificultad ---la de los responsables políticos para asumir costos inmediatos en beneficio de un resultado de largo plazo, o para reconocer públicamente errores pasados--- es universal en política. El Mercosur no es la excepción.

% ===== CAPÍTULO 8: REFLEXIÓN FINAL =====
\chapter{Discurso integracionista y práctica real}

\section{La brecha entre el discurso y la práctica}

La clase culmina con una reflexión sobre una brecha fundamental en la política latinoamericana: la distancia entre el discurso integracionista y la práctica concreta de los Estados.

El discurso de integración latinoamericana es poderoso e invoca referencias históricas profundas: los ideales bolivarianos de unidad regional, las aspiraciones sanmartinianas de liberación y unidad, la narrativa de una región que debe hablar con una voz única frente al resto del mundo.

Sin embargo, la práctica de los Estados es, frecuentemente, otra: marcada por la desconfianza mutua, por decisiones que contradicen ese discurso, por la priorización de intereses nacionales inmediatos sobre compromisos regionales.

\begin{quote}
``Hablamos permanentemente de la integración latinoamericana (...) sin embargo, en la práctica, estamos diciendo otras cosas.''
\end{quote}

\section{Indicadores de falta de voluntad política}

\subsection{Ausencia de órganos supranacionales}

Si realmente existiera una voluntad política de avanzar hacia una integración más profunda ---en el sentido de una integración comparables a la de la Unión Europea---, la estructura del Mercosur sería completamente distinta. Habría órganos independientes, con funcionarios no delegados por los gobiernos, con capacidad de decisión vinculante sobre los Estados.

La ausencia de estos órganos no responde a una imposibilidad técnica o económica. Responde a una decisión política: ningún Estado miembro está dispuesto a comprometerse de manera permanente e irrevocable con sus socios regionales en la medida que ello le quitaría control sobre sus propias políticas.

\subsection{El mecanismo de consenso como síntoma}

La persistencia del consenso como mecanismo de decisión es, en sí misma, un indicador de esa falta de voluntad. Con el consenso, cada Estado se reserva el poder de veto sobre toda decisión. Esto garantiza que ningún Estado será obligado a nada contra su voluntad en un momento dado.

Pero también significa que el bloque nunca podrá actuar con la rapidez y coherencia de una entidad verdaderamente integrada.

\section{La necesidad de discontinuidad política}

La reflexión final apunta a que, si el Mercosur realmente fuera un proyecto de integración auténtico, habría mecanismos institucionales que garantizaran su continuidad más allá de los gobiernos de turno. No es posible construir una integración profunda con gobiernos que pueden revertir, en cada cambio de gestión, lo que el anterior hizo.

La Unión Europea logró avanzar porque sus Estados miembros aceptaron ceder poder a órganos supranacionales que no cambian con cada elección. El Mercosur no ha logrado eso, porque sus Estados no están dispuestos a ceder ese poder.

\begin{important}
La pregunta fundamental no es técnica, sino política: ¿están realmente los gobiernos dispuestos a subordinar sus decisiones al proceso de integración? La respuesta actual, según la clase, es no.
\end{important}

% ===== MÉTODO CORNELL =====
\chapter*{Método Cornell de estudio}

\section*{El Mercosur como espacio de integración regional}

A continuación se presenta un cuadro de estudio según el método Cornell, que organiza los conceptos clave de esta clase en un formato que facilita el repaso activo y la recuperación de información.

\begin{table}[H]
\centering
\begin{tabularx}{\textwidth}{
    >{\raggedright\arraybackslash}p{0.28\textwidth}
    >{\raggedright\arraybackslash}X
}
\toprule
\textbf{\large Preguntas / palabras clave} &
\textbf{\large Notas / respuestas} \\
\midrule

\textbf{¿Qué es la ALADI y qué relación tiene con el Mercosur?} &
La Asociación Latinoamericana de Integración (ALADI) es el marco regional más amplio dentro del cual se permite a los Estados celebrar acuerdos de alcance parcial y conformar espacios subregionales de integración. El Mercosur es, precisamente, uno de esos espacios subregionales: no es un desprendimiento ni algo distinto de la ALADI, sino un acuerdo que se da dentro de ella. \\

\textbf{¿Quiénes fundaron el Mercosur y cuándo?} &
Argentina, Brasil, Paraguay y Uruguay, mediante el Tratado de Asunción, firmado en 1991, tras un proceso de acercamiento bilateral entre Argentina y Brasil iniciado con el restablecimiento de la democracia en la década de 1980. \\

\textbf{¿Qué es el Tratado de Asunción?} &
Es el tratado fundacional del Mercosur. Crea el Mercado Común del Sur y constituye el derecho originario del bloque, funcionando como una Constitución dentro del ordenamiento jurídico del sistema de integración. \\

\textbf{¿Qué es un mercado común y qué implica?} &
Un mercado común aspira a la libre circulación de mercaderías, servicios y personas, y a la libre circulación de los factores de producción (capital y trabajo). Esto requiere políticas uniformes entre los Estados miembros y, necesariamente, la cesión de facultades soberanas. \\

\textbf{¿Qué diferencia hay entre derecho originario y derecho derivado?} &
El derecho originario es creado directamente por los Estados Parte (el Tratado de Asunción y protocolos como Ouro Preto). El derecho derivado es el dictado por los órganos que el tratado creó, en ejercicio de las competencias que les atribuyó. La distinción depende de la \textit{función} en que actúa quien lo produce. \\

\textbf{¿Qué modificó el Protocolo de Ouro Preto?} &
Modificó la estructura orgánica del Mercosur. Aunque es un instrumento distinto del Tratado de Asunción, se incorpora a él con el mismo rango, pasando a integrar el derecho originario del bloque. \\

\textbf{¿Cuáles son los tres órganos del Mercosur y sus funciones?} &
El Consejo del Mercado Común (CMC): presidentes, ministros de Economía y cancilleres; define el rumbo político. El Grupo Mercado Común (GMC): funcionarios de Economía; ejecuta las políticas. La Comisión de Comercio (CCM): funcionarios especializados; regula cuestiones comerciales. \\

\textbf{¿Cómo se toman las decisiones en el Mercosur?} &
Por consenso, es decir, por unanimidad. Si un solo Estado miembro se opone, la norma no llega a dictarse. Esto contrasta con la Unión Europea, que usa mayorías simples, absolutas o calificadas. \\

\textbf{¿Qué efecto tiene el voto en contra de un Estado?} &
Con consenso, impide que la norma nazca jurídicamente. El Estado no queda obligado. Con mayoría, la norma se aprueba igual y el Estado disidente debe acatarla. \\

\textbf{¿Qué diferencia hay entre socio pleno y Estado asociado?} &
El socio pleno está obligado por toda la normativa del bloque y tiene derecho a voto. El Estado asociado solo se obliga por la normativa a la que adhiere expresamente y no tiene derecho a voto. \\

\textbf{¿Cuáles son los Estados socios plenos actuales?} &
Argentina, Brasil, Paraguay, Uruguay y Venezuela. Aunque formalmente Venezuela es miembro pleno, se encuentra suspendida. \\

\textbf{¿Qué es la cláusula democrática?} &
Protocolo (referido como Protocolo de Ushuaia) que exige que los Estados miembros operen dentro del Estado de Derecho y en régimen democrático. Los Estados que se apartan de estos principios pueden ser suspendidos. \\

\textbf{¿Qué pasó con Paraguay y Venezuela?} &
Paraguay fue suspendido por apartarse de premisas democráticas, precisamente cuando votaría contra el ingreso de Venezuela. Sin su voto, se aprobó Venezuela. Posteriormente, Paraguay cuestionó la validez de esa incorporación. \\

\textbf{¿Por qué es importante diferenciar entre Mercosur y Unión Europea?} &
Porque el Mercosur funciona con órganos intergubernamentales y consenso, lo que revela falta de voluntad política de avanzar hacia integración supranacional. La Unión Europea tiene órganos independientes y decide por mayoría. \\

\textbf{¿Qué muestra el caso del voto ``no positivo'' de 2008?} &
Que los gobiernos toman decisiones de política interna sin considerar sus compromisos de integración. Argentina restringía comercio internamente mientras se comprometía a libre circulación en el Mercosur. La incoherencia ejemplifica tensiones permanentes. \\

\textbf{¿Qué es una política de Estado de largo plazo?} &
Conjunto de decisiones que trascienden los ciclos electorales y los cambios de gobierno, garantizando estabilidad para que proyectos de desarrollo (industrial, de investigación, de integración) puedan consolidarse en el tiempo. \\

\textbf{¿Por qué falló la industria nacional argentina?} &
Porque decisiones de corto plazo (privatización de ferrocarriles, apertura a importaciones) se llevaron a cabo sin una política de Estado que las respaldara. Cuando se cerraron las fábricas, la capacidad productiva se perdió de manera prácticamente irreversible. \\

\textbf{¿Cuál es la reflexión final sobre integración en América Latina?} &
Existe una brecha entre el discurso integracionista (ideales bolivarianos, sanmartinianos) y la práctica real (desconfianza, priorización de intereses nacionales). La falta de órganos supranacionales no es un límite técnico, sino una decisión política de no ceder poder real. \\

\midrule

\multicolumn{2}{p{\textwidth}}{
\textbf{Síntesis:} 
El Mercosur surge como espacio subregional de integración dentro de la ALADI, fundado por Argentina, Brasil, Paraguay y Uruguay en 1991 mediante el Tratado de Asunción. Este tratado funciona como derecho originario (equivalente a una Constitución), mientras que las decisiones de sus órganos (CMC, GMC, CCM) conforman el derecho derivado. Los tres órganos funcionan por consenso, lo que refleja la falta de voluntad política de ceder poder a entidades supranacionales, contrariamente a la Unión Europea que opera por mayoría y con órganos independientes. Las tensiones constantes entre políticas internas y compromisos de integración (evidenciadas en casos como el voto ``no positivo'' o el carnet gremial de camioneros) muestran que la verdadera barrera para la integración no es institucional sino política: los gobiernos no están dispuestos a sostener políticas de largo plazo que trascienda sus ciclos electorales. La reflexión final subraya que, mientras se invoquen ideales integracionistas latinoamericanos, la práctica revela una preferencia por mantener el control nacional sobre cualquier verdadera cesión de soberanía a estructuras regionales.
} \\

\bottomrule
\end{tabularx}
\end{table}

% ===== FIN DEL DOCUMENTO =====
\end{document}
